% ============================================================
% 图表 LaTeX 引用文件 — 自动生成
% 在论文主文件中使用 % ============================================================
% 图表 LaTeX 引用文件 — 自动生成
% 在论文主文件中使用 % ============================================================
% 图表 LaTeX 引用文件 — 自动生成
% 在论文主文件中使用 % ============================================================
% 图表 LaTeX 引用文件 — 自动生成
% 在论文主文件中使用 \input{figures/latex_includes.tex} 引入
% ============================================================

% --- 问题一 ---

\begin{figure}[H]
\centering
\includegraphics[width=0.95\textwidth]{figures/fig_q1_power_balance.pdf}
\caption{典型日功率平衡曲线（风电、光伏出力与用电负荷、购售电功率随时间变化）}
\label{fig:q1_power_balance}
\end{figure}

\begin{figure}[H]
\centering
\includegraphics[width=0.75\textwidth]{figures/fig_q1_indicators.pdf}
\caption{绿电直连三项指标实际值与阈值要求对比}
\label{fig:q1_indicators}
\end{figure}

% --- 问题二 ---

\begin{figure}[H]
\centering
\includegraphics[width=0.8\textwidth]{figures/fig_q2_cost_vs_production.pdf}
\caption{不同日产氨量下的吨氨成本变化趋势}
\label{fig:q2_cost_vs_production}
\end{figure}

\begin{figure}[H]
\centering
\includegraphics[width=0.95\textwidth]{figures/fig_q2_schedule_gantt.pdf}
\caption{各产量水平下最优24小时开停机时段安排}
\label{fig:q2_schedule_gantt}
\end{figure}

\begin{figure}[H]
\centering
\includegraphics[width=0.7\textwidth]{figures/fig_q2_scenarios_heatmap.pdf}
\caption{24种风光场景下吨氨成本热力图（6种风电$\times$4种光伏出力水平）}
\label{fig:q2_scenarios_heatmap}
\end{figure}

\begin{figure}[H]
\centering
\includegraphics[width=0.9\textwidth]{figures/fig_q2_annual_cost.pdf}
\caption{全年吨氨成本分布曲线（问题二离散调节与问题三连续调节对比）}
\label{fig:q2_annual_cost}
\end{figure}

\begin{figure}[H]
\centering
\includegraphics[width=0.7\textwidth]{figures/fig_q2_green_indicators.pdf}
\caption{全年绿电直连指标满足情况统计（离散调节与连续调节对比）}
\label{fig:q2_green_indicators}
\end{figure}

% --- 问题三 ---

\begin{figure}[H]
\centering
\includegraphics[width=0.95\textwidth]{figures/fig_q3_dispatch.pdf}
\caption{典型场景下连续调度功率曲线与制氢氨负荷率}
\label{fig:q3_dispatch}
\end{figure}

\begin{figure}[H]
\centering
\includegraphics[width=0.9\textwidth]{figures/fig_q3_annual_cost.pdf}
\caption{全年吨氨成本排序分布曲线（问题三与问题二对比）}
\label{fig:q3_annual_cost}
\end{figure}

\begin{figure}[H]
\centering
\includegraphics[width=0.8\textwidth]{figures/fig_q3_comparison.pdf}
\caption{连续调节相对离散调节各项指标变化率}
\label{fig:q3_comparison}
\end{figure}

% --- 问题四 ---

\begin{figure}[H]
\centering
\includegraphics[width=0.95\textwidth]{figures/fig_q4_offgrid_production.pdf}
\caption{24种风光场景下离网制氨日产量对比（有/无储能）}
\label{fig:q4_offgrid_production}
\end{figure}

\begin{figure}[H]
\centering
\includegraphics[width=0.95\textwidth]{figures/fig_q4_storage_dispatch.pdf}
\caption{最大弃电场景下储能参与的功率调度曲线与SOC变化}
\label{fig:q4_storage_dispatch}
\end{figure}

\begin{figure}[H]
\centering
\includegraphics[width=0.85\textwidth]{figures/fig_q4_economics.pdf}
\caption{离网与联网运行模式吨氨成本分解对比（瀑布图）}
\label{fig:q4_economics}
\end{figure}

% --- 问题五 ---

\begin{figure}[H]
\centering
\includegraphics[width=0.7\textwidth]{figures/fig_q5_radar.pdf}
\caption{绿电园区容量渗透率对电力系统影响多维评价}
\label{fig:q5_radar}
\end{figure}

% === DrawIO/TikZ 架构图与流程图 ===

% 技术路线图
\begin{figure}[H]
\centering
\includegraphics[width=\textwidth,height=0.45\textheight,keepaspectratio]{figures/fig_roadmap.pdf}
\caption{绿电直连型电氢氨园区优化运行整体技术路线图}\label{fig:roadmap}
\end{figure}

% 园区能量流向架构图
\begin{figure}[H]
\centering
\includegraphics[width=0.9\textwidth,height=0.3\textheight,keepaspectratio]{figures/fig_energy_flow.pdf}
\caption{绿电直连型电氢氨园区能量流向架构图}\label{fig:energy_flow}
\end{figure}

% 问题一求解流程图
\begin{figure}[H]
\centering
\includegraphics[width=0.85\textwidth,height=0.38\textheight,keepaspectratio]{figures/fig_flow_q1.pdf}
\caption{问题一求解流程图（典型日运行指标分析）}\label{fig:flow_q1}
\end{figure}

% 问题二求解流程图
\begin{figure}[H]
\centering
\includegraphics[width=0.85\textwidth,height=0.38\textheight,keepaspectratio]{figures/fig_flow_q2.pdf}
\caption{问题二求解流程图（离散制氨调节优化）}\label{fig:flow_q2}
\end{figure}

% 问题三求解流程图
\begin{figure}[H]
\centering
\includegraphics[width=0.85\textwidth,height=0.38\textheight,keepaspectratio]{figures/fig_flow_q3.pdf}
\caption{问题三求解流程图（连续功率调度优化）}\label{fig:flow_q3}
\end{figure}

% 问题四求解流程图
\begin{figure}[H]
\centering
\includegraphics[width=0.85\textwidth,height=0.38\textheight,keepaspectratio]{figures/fig_flow_q4.pdf}
\caption{问题四求解流程图（离网运行与储能配置）}\label{fig:flow_q4}
\end{figure}

% 储能状态转移方程示意图 (TikZ)
\begin{figure}[H]
\centering
\includegraphics[width=0.85\textwidth,height=0.35\textheight,keepaspectratio]{figures/tikz_storage_model.pdf}
\caption{储能系统SOC状态转移方程示意图}\label{fig:tikz_storage}
\end{figure}
 引入
% ============================================================

% --- 问题一 ---

\begin{figure}[H]
\centering
\includegraphics[width=0.95\textwidth]{figures/fig_q1_power_balance.pdf}
\caption{典型日功率平衡曲线（风电、光伏出力与用电负荷、购售电功率随时间变化）}
\label{fig:q1_power_balance}
\end{figure}

\begin{figure}[H]
\centering
\includegraphics[width=0.75\textwidth]{figures/fig_q1_indicators.pdf}
\caption{绿电直连三项指标实际值与阈值要求对比}
\label{fig:q1_indicators}
\end{figure}

% --- 问题二 ---

\begin{figure}[H]
\centering
\includegraphics[width=0.8\textwidth]{figures/fig_q2_cost_vs_production.pdf}
\caption{不同日产氨量下的吨氨成本变化趋势}
\label{fig:q2_cost_vs_production}
\end{figure}

\begin{figure}[H]
\centering
\includegraphics[width=0.95\textwidth]{figures/fig_q2_schedule_gantt.pdf}
\caption{各产量水平下最优24小时开停机时段安排}
\label{fig:q2_schedule_gantt}
\end{figure}

\begin{figure}[H]
\centering
\includegraphics[width=0.7\textwidth]{figures/fig_q2_scenarios_heatmap.pdf}
\caption{24种风光场景下吨氨成本热力图（6种风电$\times$4种光伏出力水平）}
\label{fig:q2_scenarios_heatmap}
\end{figure}

\begin{figure}[H]
\centering
\includegraphics[width=0.9\textwidth]{figures/fig_q2_annual_cost.pdf}
\caption{全年吨氨成本分布曲线（问题二离散调节与问题三连续调节对比）}
\label{fig:q2_annual_cost}
\end{figure}

\begin{figure}[H]
\centering
\includegraphics[width=0.7\textwidth]{figures/fig_q2_green_indicators.pdf}
\caption{全年绿电直连指标满足情况统计（离散调节与连续调节对比）}
\label{fig:q2_green_indicators}
\end{figure}

% --- 问题三 ---

\begin{figure}[H]
\centering
\includegraphics[width=0.95\textwidth]{figures/fig_q3_dispatch.pdf}
\caption{典型场景下连续调度功率曲线与制氢氨负荷率}
\label{fig:q3_dispatch}
\end{figure}

\begin{figure}[H]
\centering
\includegraphics[width=0.9\textwidth]{figures/fig_q3_annual_cost.pdf}
\caption{全年吨氨成本排序分布曲线（问题三与问题二对比）}
\label{fig:q3_annual_cost}
\end{figure}

\begin{figure}[H]
\centering
\includegraphics[width=0.8\textwidth]{figures/fig_q3_comparison.pdf}
\caption{连续调节相对离散调节各项指标变化率}
\label{fig:q3_comparison}
\end{figure}

% --- 问题四 ---

\begin{figure}[H]
\centering
\includegraphics[width=0.95\textwidth]{figures/fig_q4_offgrid_production.pdf}
\caption{24种风光场景下离网制氨日产量对比（有/无储能）}
\label{fig:q4_offgrid_production}
\end{figure}

\begin{figure}[H]
\centering
\includegraphics[width=0.95\textwidth]{figures/fig_q4_storage_dispatch.pdf}
\caption{最大弃电场景下储能参与的功率调度曲线与SOC变化}
\label{fig:q4_storage_dispatch}
\end{figure}

\begin{figure}[H]
\centering
\includegraphics[width=0.85\textwidth]{figures/fig_q4_economics.pdf}
\caption{离网与联网运行模式吨氨成本分解对比（瀑布图）}
\label{fig:q4_economics}
\end{figure}

% --- 问题五 ---

\begin{figure}[H]
\centering
\includegraphics[width=0.7\textwidth]{figures/fig_q5_radar.pdf}
\caption{绿电园区容量渗透率对电力系统影响多维评价}
\label{fig:q5_radar}
\end{figure}

% === DrawIO/TikZ 架构图与流程图 ===

% 技术路线图
\begin{figure}[H]
\centering
\includegraphics[width=\textwidth,height=0.45\textheight,keepaspectratio]{figures/fig_roadmap.pdf}
\caption{绿电直连型电氢氨园区优化运行整体技术路线图}\label{fig:roadmap}
\end{figure}

% 园区能量流向架构图
\begin{figure}[H]
\centering
\includegraphics[width=0.9\textwidth,height=0.3\textheight,keepaspectratio]{figures/fig_energy_flow.pdf}
\caption{绿电直连型电氢氨园区能量流向架构图}\label{fig:energy_flow}
\end{figure}

% 问题一求解流程图
\begin{figure}[H]
\centering
\includegraphics[width=0.85\textwidth,height=0.38\textheight,keepaspectratio]{figures/fig_flow_q1.pdf}
\caption{问题一求解流程图（典型日运行指标分析）}\label{fig:flow_q1}
\end{figure}

% 问题二求解流程图
\begin{figure}[H]
\centering
\includegraphics[width=0.85\textwidth,height=0.38\textheight,keepaspectratio]{figures/fig_flow_q2.pdf}
\caption{问题二求解流程图（离散制氨调节优化）}\label{fig:flow_q2}
\end{figure}

% 问题三求解流程图
\begin{figure}[H]
\centering
\includegraphics[width=0.85\textwidth,height=0.38\textheight,keepaspectratio]{figures/fig_flow_q3.pdf}
\caption{问题三求解流程图（连续功率调度优化）}\label{fig:flow_q3}
\end{figure}

% 问题四求解流程图
\begin{figure}[H]
\centering
\includegraphics[width=0.85\textwidth,height=0.38\textheight,keepaspectratio]{figures/fig_flow_q4.pdf}
\caption{问题四求解流程图（离网运行与储能配置）}\label{fig:flow_q4}
\end{figure}

% 储能状态转移方程示意图 (TikZ)
\begin{figure}[H]
\centering
\includegraphics[width=0.85\textwidth,height=0.35\textheight,keepaspectratio]{figures/tikz_storage_model.pdf}
\caption{储能系统SOC状态转移方程示意图}\label{fig:tikz_storage}
\end{figure}
 引入
% ============================================================

% --- 问题一 ---

\begin{figure}[H]
\centering
\includegraphics[width=0.95\textwidth]{figures/fig_q1_power_balance.pdf}
\caption{典型日功率平衡曲线（风电、光伏出力与用电负荷、购售电功率随时间变化）}
\label{fig:q1_power_balance}
\end{figure}

\begin{figure}[H]
\centering
\includegraphics[width=0.75\textwidth]{figures/fig_q1_indicators.pdf}
\caption{绿电直连三项指标实际值与阈值要求对比}
\label{fig:q1_indicators}
\end{figure}

% --- 问题二 ---

\begin{figure}[H]
\centering
\includegraphics[width=0.8\textwidth]{figures/fig_q2_cost_vs_production.pdf}
\caption{不同日产氨量下的吨氨成本变化趋势}
\label{fig:q2_cost_vs_production}
\end{figure}

\begin{figure}[H]
\centering
\includegraphics[width=0.95\textwidth]{figures/fig_q2_schedule_gantt.pdf}
\caption{各产量水平下最优24小时开停机时段安排}
\label{fig:q2_schedule_gantt}
\end{figure}

\begin{figure}[H]
\centering
\includegraphics[width=0.7\textwidth]{figures/fig_q2_scenarios_heatmap.pdf}
\caption{24种风光场景下吨氨成本热力图（6种风电$\times$4种光伏出力水平）}
\label{fig:q2_scenarios_heatmap}
\end{figure}

\begin{figure}[H]
\centering
\includegraphics[width=0.9\textwidth]{figures/fig_q2_annual_cost.pdf}
\caption{全年吨氨成本分布曲线（问题二离散调节与问题三连续调节对比）}
\label{fig:q2_annual_cost}
\end{figure}

\begin{figure}[H]
\centering
\includegraphics[width=0.7\textwidth]{figures/fig_q2_green_indicators.pdf}
\caption{全年绿电直连指标满足情况统计（离散调节与连续调节对比）}
\label{fig:q2_green_indicators}
\end{figure}

% --- 问题三 ---

\begin{figure}[H]
\centering
\includegraphics[width=0.95\textwidth]{figures/fig_q3_dispatch.pdf}
\caption{典型场景下连续调度功率曲线与制氢氨负荷率}
\label{fig:q3_dispatch}
\end{figure}

\begin{figure}[H]
\centering
\includegraphics[width=0.9\textwidth]{figures/fig_q3_annual_cost.pdf}
\caption{全年吨氨成本排序分布曲线（问题三与问题二对比）}
\label{fig:q3_annual_cost}
\end{figure}

\begin{figure}[H]
\centering
\includegraphics[width=0.8\textwidth]{figures/fig_q3_comparison.pdf}
\caption{连续调节相对离散调节各项指标变化率}
\label{fig:q3_comparison}
\end{figure}

% --- 问题四 ---

\begin{figure}[H]
\centering
\includegraphics[width=0.95\textwidth]{figures/fig_q4_offgrid_production.pdf}
\caption{24种风光场景下离网制氨日产量对比（有/无储能）}
\label{fig:q4_offgrid_production}
\end{figure}

\begin{figure}[H]
\centering
\includegraphics[width=0.95\textwidth]{figures/fig_q4_storage_dispatch.pdf}
\caption{最大弃电场景下储能参与的功率调度曲线与SOC变化}
\label{fig:q4_storage_dispatch}
\end{figure}

\begin{figure}[H]
\centering
\includegraphics[width=0.85\textwidth]{figures/fig_q4_economics.pdf}
\caption{离网与联网运行模式吨氨成本分解对比（瀑布图）}
\label{fig:q4_economics}
\end{figure}

% --- 问题五 ---

\begin{figure}[H]
\centering
\includegraphics[width=0.7\textwidth]{figures/fig_q5_radar.pdf}
\caption{绿电园区容量渗透率对电力系统影响多维评价}
\label{fig:q5_radar}
\end{figure}

% === DrawIO/TikZ 架构图与流程图 ===

% 技术路线图
\begin{figure}[H]
\centering
\includegraphics[width=\textwidth,height=0.45\textheight,keepaspectratio]{figures/fig_roadmap.pdf}
\caption{绿电直连型电氢氨园区优化运行整体技术路线图}\label{fig:roadmap}
\end{figure}

% 园区能量流向架构图
\begin{figure}[H]
\centering
\includegraphics[width=0.9\textwidth,height=0.3\textheight,keepaspectratio]{figures/fig_energy_flow.pdf}
\caption{绿电直连型电氢氨园区能量流向架构图}\label{fig:energy_flow}
\end{figure}

% 问题一求解流程图
\begin{figure}[H]
\centering
\includegraphics[width=0.85\textwidth,height=0.38\textheight,keepaspectratio]{figures/fig_flow_q1.pdf}
\caption{问题一求解流程图（典型日运行指标分析）}\label{fig:flow_q1}
\end{figure}

% 问题二求解流程图
\begin{figure}[H]
\centering
\includegraphics[width=0.85\textwidth,height=0.38\textheight,keepaspectratio]{figures/fig_flow_q2.pdf}
\caption{问题二求解流程图（离散制氨调节优化）}\label{fig:flow_q2}
\end{figure}

% 问题三求解流程图
\begin{figure}[H]
\centering
\includegraphics[width=0.85\textwidth,height=0.38\textheight,keepaspectratio]{figures/fig_flow_q3.pdf}
\caption{问题三求解流程图（连续功率调度优化）}\label{fig:flow_q3}
\end{figure}

% 问题四求解流程图
\begin{figure}[H]
\centering
\includegraphics[width=0.85\textwidth,height=0.38\textheight,keepaspectratio]{figures/fig_flow_q4.pdf}
\caption{问题四求解流程图（离网运行与储能配置）}\label{fig:flow_q4}
\end{figure}

% 储能状态转移方程示意图 (TikZ)
\begin{figure}[H]
\centering
\includegraphics[width=0.85\textwidth,height=0.35\textheight,keepaspectratio]{figures/tikz_storage_model.pdf}
\caption{储能系统SOC状态转移方程示意图}\label{fig:tikz_storage}
\end{figure}
 引入
% ============================================================

% --- 问题一 ---

\begin{figure}[H]
\centering
\includegraphics[width=0.95\textwidth]{figures/fig_q1_power_balance.pdf}
\caption{典型日功率平衡曲线（风电、光伏出力与用电负荷、购售电功率随时间变化）}
\label{fig:q1_power_balance}
\end{figure}

\begin{figure}[H]
\centering
\includegraphics[width=0.75\textwidth]{figures/fig_q1_indicators.pdf}
\caption{绿电直连三项指标实际值与阈值要求对比}
\label{fig:q1_indicators}
\end{figure}

% --- 问题二 ---

\begin{figure}[H]
\centering
\includegraphics[width=0.8\textwidth]{figures/fig_q2_cost_vs_production.pdf}
\caption{不同日产氨量下的吨氨成本变化趋势}
\label{fig:q2_cost_vs_production}
\end{figure}

\begin{figure}[H]
\centering
\includegraphics[width=0.95\textwidth]{figures/fig_q2_schedule_gantt.pdf}
\caption{各产量水平下最优24小时开停机时段安排}
\label{fig:q2_schedule_gantt}
\end{figure}

\begin{figure}[H]
\centering
\includegraphics[width=0.7\textwidth]{figures/fig_q2_scenarios_heatmap.pdf}
\caption{24种风光场景下吨氨成本热力图（6种风电$\times$4种光伏出力水平）}
\label{fig:q2_scenarios_heatmap}
\end{figure}

\begin{figure}[H]
\centering
\includegraphics[width=0.9\textwidth]{figures/fig_q2_annual_cost.pdf}
\caption{全年吨氨成本分布曲线（问题二离散调节与问题三连续调节对比）}
\label{fig:q2_annual_cost}
\end{figure}

\begin{figure}[H]
\centering
\includegraphics[width=0.7\textwidth]{figures/fig_q2_green_indicators.pdf}
\caption{全年绿电直连指标满足情况统计（离散调节与连续调节对比）}
\label{fig:q2_green_indicators}
\end{figure}

% --- 问题三 ---

\begin{figure}[H]
\centering
\includegraphics[width=0.95\textwidth]{figures/fig_q3_dispatch.pdf}
\caption{典型场景下连续调度功率曲线与制氢氨负荷率}
\label{fig:q3_dispatch}
\end{figure}

\begin{figure}[H]
\centering
\includegraphics[width=0.9\textwidth]{figures/fig_q3_annual_cost.pdf}
\caption{全年吨氨成本排序分布曲线（问题三与问题二对比）}
\label{fig:q3_annual_cost}
\end{figure}

\begin{figure}[H]
\centering
\includegraphics[width=0.8\textwidth]{figures/fig_q3_comparison.pdf}
\caption{连续调节相对离散调节各项指标变化率}
\label{fig:q3_comparison}
\end{figure}

% --- 问题四 ---

\begin{figure}[H]
\centering
\includegraphics[width=0.95\textwidth]{figures/fig_q4_offgrid_production.pdf}
\caption{24种风光场景下离网制氨日产量对比（有/无储能）}
\label{fig:q4_offgrid_production}
\end{figure}

\begin{figure}[H]
\centering
\includegraphics[width=0.95\textwidth]{figures/fig_q4_storage_dispatch.pdf}
\caption{最大弃电场景下储能参与的功率调度曲线与SOC变化}
\label{fig:q4_storage_dispatch}
\end{figure}

\begin{figure}[H]
\centering
\includegraphics[width=0.85\textwidth]{figures/fig_q4_economics.pdf}
\caption{离网与联网运行模式吨氨成本分解对比（瀑布图）}
\label{fig:q4_economics}
\end{figure}

% --- 问题五 ---

\begin{figure}[H]
\centering
\includegraphics[width=0.7\textwidth]{figures/fig_q5_radar.pdf}
\caption{绿电园区容量渗透率对电力系统影响多维评价}
\label{fig:q5_radar}
\end{figure}

% === DrawIO/TikZ 架构图与流程图 ===

% 技术路线图
\begin{figure}[H]
\centering
\includegraphics[width=\textwidth,height=0.45\textheight,keepaspectratio]{figures/fig_roadmap.pdf}
\caption{绿电直连型电氢氨园区优化运行整体技术路线图}\label{fig:roadmap}
\end{figure}

% 园区能量流向架构图
\begin{figure}[H]
\centering
\includegraphics[width=0.9\textwidth,height=0.3\textheight,keepaspectratio]{figures/fig_energy_flow.pdf}
\caption{绿电直连型电氢氨园区能量流向架构图}\label{fig:energy_flow}
\end{figure}

% 问题一求解流程图
\begin{figure}[H]
\centering
\includegraphics[width=0.85\textwidth,height=0.38\textheight,keepaspectratio]{figures/fig_flow_q1.pdf}
\caption{问题一求解流程图（典型日运行指标分析）}\label{fig:flow_q1}
\end{figure}

% 问题二求解流程图
\begin{figure}[H]
\centering
\includegraphics[width=0.85\textwidth,height=0.38\textheight,keepaspectratio]{figures/fig_flow_q2.pdf}
\caption{问题二求解流程图（离散制氨调节优化）}\label{fig:flow_q2}
\end{figure}

% 问题三求解流程图
\begin{figure}[H]
\centering
\includegraphics[width=0.85\textwidth,height=0.38\textheight,keepaspectratio]{figures/fig_flow_q3.pdf}
\caption{问题三求解流程图（连续功率调度优化）}\label{fig:flow_q3}
\end{figure}

% 问题四求解流程图
\begin{figure}[H]
\centering
\includegraphics[width=0.85\textwidth,height=0.38\textheight,keepaspectratio]{figures/fig_flow_q4.pdf}
\caption{问题四求解流程图（离网运行与储能配置）}\label{fig:flow_q4}
\end{figure}

% 储能状态转移方程示意图 (TikZ)
\begin{figure}[H]
\centering
\includegraphics[width=0.85\textwidth,height=0.35\textheight,keepaspectratio]{figures/tikz_storage_model.pdf}
\caption{储能系统SOC状态转移方程示意图}\label{fig:tikz_storage}
\end{figure}
